% page 380 of the facsimile (image p-379 of the scan)
% block 149.9 x 237.7 mm ; left margin 8.0 mm
\pgc{-12.700mm}{0.000mm}{
\l{\cel{3}372}
\l{}
\l{mikra. I. Qua ne atingas la dimensioni ordinara, partiku-}
\l{\cel{3}lare pri longeso, alteso, qualeso o quanteso. -II. Qua}
\l{\cel{3}ne atingas la nivelo ordinara pri rango en la hierarkio,}
\l{\cel{3}merito, la qualesi di kordio o di mento. - DEFIRS.}
\l{}
\l{\sou{mikr}o-\cel{3}Prefixo ciencala : la milionimo de...}
\l{}
\l{mikrobo. Celulo vivanta, di vejetanto od animalo, qua apar-}
\l{\cel{3}tenas a la mondo di la infinite mikra, faktoro di ula}
\l{\cel{3}fermentaci od ula morbi infektiva. - DEFIRS.}
\l{}
\l{\sou{mikrocentro}. (\sou{biol.}) Parto centra di la sfero atraktiva,}
\l{\cel{3}konsistanta ye plura kromosomi.}
\l{}
\l{\sou{mikrofono}. (\sou{fiziko}).Aparato akustikala, destinita ad am-}
\l{\cel{3}pligar la soni. - DEFIRS.}
\l{}
\l{\sou{mikrogameto}. (\sou{biol.}) Elemento genetiva, mikra e moviva, kun}
\l{\cel{3}karakteri maskulala.}
\l{}
\l{\sou{mikrograf}o. Sorto di pantografo, qua posibligas desegnar}
\l{\cel{3}figuri extreme mikra.}
\l{}
\l{\sou{mikrografio}. La cienco qua traktas la preparo di la objek-}
\l{\cel{3}ti por lia observado per mikroskopo, e la deskripto di}
\l{\cel{3}olci pos lia exameneso per ta instrumento. - DEFIRS.}
\l{}
\l{\sou{mikrometro}. (\sou{metrologio)}. Instrumento destinita a la mezuro}
\l{\cel{3}di mikra objekti e mikra imaji. - DEFIRS.}
\l{}
\l{\sou{mikrono}. (\sou{teknol}.)La milima parto di la milimetro - unajo}
\l{\cel{3}di mezuro qua adoptesis por la observi mikroskopala. -}
\l{\cel{3}DEFIRS.}
\l{}
\l{\sou{mikronukleo}. (\sou{biol.})\cel{2}La nukleo maxim mikra che la infuzo-}
\l{\cel{3}rii, opoze a la grosa nukleo (\sou{makronukleo}).}
\l{}
\l{\sou{mikropilo}. (\sou{biol.}) Mikra orifico di la membrano vitelala}
\l{\cel{3}di la ovulo, tra qua pasas (konjektate) la spermatozoido}
\l{\cel{3}lor la fekundigo. -}
\l{}
\l{\sou{mikroskopo}. Instrumento optikala, ek lensi dispozita tale}
\l{\cel{3}ke li igas la mikra objekti semblar plu grosa kam se}
\l{\cel{3}regardata per okulo nuda. - DEFIRS.}
\l{}
\l{\sou{mikrosomo}. (\sou{biol}.) (\sou{histologio}).Granularo tenua quan onu}
\l{\cel{3}renkontras en la parto figurala di citoplasmo. DEFIS.}
\l{}
\l{\sou{mikrosporo}. Nomo di ta spori maxim mikra di la vejetanti,}
\l{\cel{3}qui genitas la protali maskula. - DEFIS.}
\l{}
\l{\sou{mikrosporangio}. Organo, che algi diversa e kriptogami}
\l{\cel{3}vaskuloza, en qua produktesas la mikrospori. - DEFIS.}
}
